%-------------------------------------------------------------------------------
%	SECTION TITLE
%-------------------------------------------------------------------------------
\cvsection{Formación académica y posdoctoral}


%-------------------------------------------------------------------------------
%	CONTENT
%-------------------------------------------------------------------------------
\begin{cventries}

%---------------------------------------------------------
%  \cventry
%    {B.S. in Computer Science and Engineering} % Degree
%    {POSTECH(Pohang University of Science and Technology)} % Institution
%    {Pohang, S.Korea} % Location
%    {Mar. 2010 - Aug. 2017} % Date(s)
%    {
%      \begin{cvitems} % Description(s) bullet points
%        \item {Got a Chun Shin-Il Scholarship which is given to promising students in CSE Dept.}
%      \end{cvitems}
%    }



\cventry
    {CONICET y UPC}
    {Beca posdoctoral}
    {Argentina/España}
    {2015 - 2018}
    {
      \begin{cvitems} % Description(s) bullet points
        \item {Tema de investigación: evolución térmica de la litosfera oceánica.
Directores: Dr. Cristallini (CONICET-UBA) y Dr. Zlotnik (UPC).}
      \end{cvitems}
    }

\cventry
    {Stanford University}
    {Investigador visitante en Stanford}
    {Estados Unidos}
    {2013}
    {
      \begin{cvitems} % Description(s) bullet points
        \item {\emph{Cursos realizados: Geología Estructural y MATLAB aplicado a Geociencias.}}
      \end{cvitems}
    }

\cventry
    {Universidad de Buenos Aires}
    {Doctorado en Ciencias Geológicas}
    {Argentina}
    {2010 - 2015}
    {
      \begin{cvitems} % Description(s) bullet points
        \item {Tesis: Predicción del fracturamiento en superficies geológicas plegadas mediante modelado numérico. Director: Dr. Cristallini.}
      \end{cvitems}
    }

\cventry
    {Universidad de Buenos Aires}
    {Licenciatura en Ciencias Geológicas}
    {Argentina}
    {2005 - 2010}
    {
      \begin{cvitems} % Description(s) bullet points
        \item {Tesis: Estructura y geología del anticlinal compuesto Tres Cruces, Jujuy, Argentina. Director: Dr. Cristallini.}
      \end{cvitems}
    }

%---------------------------------------------------------
\end{cventries}
